\begin{lemma}
\label{thm: get epsilon close many times}
    If $\best{q}$ is globally asymptotically stable for the deterministic difference inclusion~\eqref{eq: deterministic mces in stochastic section} with any sequences $(\lrate_k)$ and $(\updateset_k)$ that satisfy \autoref{asp: stochastic approximation step size},
    % under \autoref{asp: stochastic approximation step size}, 
    then 
    % the stochastic algorithm \eqref{eq: difference inclusion to show robustness noisy} gets arbitrarily close to $q_\opt$ at least once.
    % In other words, 
    for any $\epsilon>0$, there exists a time step $K>0$ when the solution generated by MC-OPI satisfies 
    % for a sequence $(\qval_\kpz)$ generated by algorithm \eqref{eq: difference inclusion to show robustness noisy}, it holds that    
    \begin{align}
        \qval_K \in \best{q} + \epsilon \ball
        % \forall \epsilon>0, \exists K >0 : \qval_K \in \qval_\opt + \epsilon \Bar{\ball}
    \end{align}
    \textcolor{red}{there exist an infinite number of time steps}
\end{lemma}
\begin{proof}
    % Since the difference inclusion \eqref{eq: difference inclusion to show robustness noisy} is outer semicontinuous, non-empty and compact for every $q \in \qset$,
    % and that, by assumption, it is globally asymptotically stable, 
    % \citet[Theorem~12]{Kellett2004SmoothInclusions} and \citet[Theorem~13.1]{Teel2017RecentInclusions} 
    % show that its convergence is robust to perturbations on $q$ and $i$.
    % 
    Example~2 in \citep{Kellett2007SufficientInclusions} shows that the difference inclusion \eqref{eq: deterministic mces in stochastic section} is robust to small perturbations because it is globally asymptotically stable and outer semicontinuous for every time step $k \in \integers_{\geq 0}$.
    % 
    In particular, there exists a positive definite function $\sigma : \qset \to \real_{\geq 0}$ satisfying
    \begin{align*}
        \sigma(\qval)>0 \quad \forall \ q \neq \qval\opt ,
    \end{align*}
    such that $\qval\opt$ is globally asymptotically stable for the perturbed inclusion
    \begin{align}
    % \label{eq: perturbed inclusion}
        q_\kpo
        \in
        \big\{ \ \qval_\kpz + \lrate_\kpz \update \big( q_{\pol} - \qval_\kpz \big) + \sigma(q_k) \Bar{\ball}
        \ \text{ with } \
        \update \in \updateset_k(q_k),
        q_{\pol} \in Q(q_k)
        \ \big\} .
    \end{align}
    % \begin{align}
    % \label{eq: perturbed inclusion}
    %     \begin{bmatrix}
    %         q_\kpo \\
    %         i_\kpo
    %     \end{bmatrix}
    %     \in
    %     \Bigg\{
    %     \begin{bmatrix}
    %         \qval_k + a(i_k) \update_k \big( q_{\pol_k} - \qval_k \big) + \sigma(q_k) \Bar{\ball}
    %         \\
    %         i_k + 1
    %     \end{bmatrix}
    %     \text{ with }
    %     \update_k \in \updateset(q_k),
    %     q_{\pol_k} \in Q(q_k)
    %     \Bigg\}
    % \end{align}
    % \begin{subequations}
    % \label{eq: perturbed inclusion}
    % \begin{align}
    %     \begin{bmatrix}
    %         q_+ \\
    %         k_+
    %     \end{bmatrix}
    %     &\in
    %     \begin{bmatrix}
    %         \qval + \lrate(\kpz) \update \big( \qval_{\pol} - \qval \big) + \sigma(q) \Bar{\ball}
    %         % (1 - \lrate(\kpz) \update) ( \qval + \sigma(q) \Bar{\ball}) + \lrate(\kpz) \update \qval_{\pol} + \sigma(q) \Bar{\ball}
    %         \\
    %         k+1
    %     \end{bmatrix}
    %     % \\
    %     % \update &\in M(q + \sigma(q) \Bar{\ball})
    %     % \\
    %     % \qval_\pol &\in Q(q + \sigma(q) \Bar{\ball})
    %     \\
    %     \update &\in M(q)
    %     \\
    %     \qval_\pol &\in Q(q)
    % \end{align}
    % \end{subequations}
    Then, for any $\epsilon > 0$, we can construct a region
    \begin{align}
        \qset_\epsilon
        \coloneqq
        \big\{ \ q \in \qset \mid \worst{q} \leq q + \epsilon \ball \leq \best{q} \text{~and~} q \notin \best{q} + \epsilon \ball \ \big\}
    \end{align}
    and the associated perturbation radius
    \begin{align}
        \sigma_\epsilon \coloneqq
        \min_{\qval \in \qset_\epsilon} \sigma(q) > 0 .
    \end{align}
    % Note that the minimum is well defined because $\qset_\epsilon$ is closed and $\sigma$ is lower-bounded.
    Since the sequence $(q_\kpz)$ generated by MC-OPI converges to $\qset_\epsilon$ (\autoref{thm: bounded solution}) 
    and the sequence $(\lrate_k \update_\kpz w_\kpz)$ converges to zero with probability one (\autoref{thm: effects of noise decrease}),
    there exists a time step $K$ such that for all $k>K$
    \begin{align}
        \max_{s,a} | \lrate_k \update_k w_k(s,a) | < \sigma_\epsilon .
    \end{align}
    As a result, after time step $K$, MC-OPI behaves like the difference inclusion~\eqref{eq: perturbed inclusion} in region $\qset_\epsilon$, and eventually enters the $\epsilon$-ball centered in $\best{q}$.

    Besides, every time MC-OPI exits this $\epsilon$-ball, it enters $\qset_\epsilon$, behaves like the difference inclusion~\eqref{eq: perturbed inclusion} again, and eventually enters the $\epsilon$-ball centered in $\best{q}$.
\end{proof}